\section{Static sign certificates with all training-gate changes}
\label{sec:static}
This section concerns parameter neighborhoods, not their reachability from
initialization. Its dataset consists of the exact binary numbers in the
accompanying \texttt{canonical\_witness\_states.npz}. Rounding a Gaussian
sample is not an iid Gaussian theorem.

Let $y$ be a parameter center, $R>0$, $P=\|y\|$, $\bar P=P+R$, and
$D=PR+R^2/2$. A cluster index set $C$ uses the full-dataset normalization
$\|z\|_C^2=N^{-1}\sum_{n\in C}\|z_n\|^2$. Put
$\lambda_C\ge\|N^{-1}\sum_{n\in C}x_nx_n^T\|$.
Write $F_C$ for that cluster's contribution to the full gradient field,
using the \emph{full-network residual}; $F=F_1+F_2$.
Let $F_C^0$ denote its polynomial extension with the training gates fixed
at $y$. No fixed-gate hypothesis is imposed on the actual field in the ball.

For either potential used here, set $g=\nabla\Phi$, $G_0=\|g(y)\|$.
For the empirical statistic use $\Phi=\Phi_I$ from
Proposition~\ref{prop:potential}. For the full derivative, fix the actual
probe-active set at $y$ and use
$\Phi(\theta)=\frac12\|M_{\rm act}(\theta)\xi-\xi\|^2$.
This second choice is valid throughout the ball if
$\min_i|u_i(y)^T\xi|>R\|\xi\|$.
In either case let $v$ be its input vector, $I$ its fixed index set,
$P_I=\|(u_i,w_i)_{i\in I}(y)\|$, and $\ell_y$ its output gradient.
Define
\begin{align*}
 L_C&=\lambda_C\bar P^2+
 \sqrt{\lambda_C}\big(\|e(y)\|_C+\sqrt{\lambda_C}D\big),\\
 H_g&=\|v\|^2(P_I+R)^2+
 \|v\|\big(\|\ell_y\|+\|v\|(P_IR+R^2/2)\big).
\end{align*}

Here are explicit bounds for the training-gate correction. Put
$q_{ni}=u_i(y)^Tx_n$, $s_n=\|x_n\|$, $e_n=e(y;x_n)$, and
\begin{align*}
 d_{ni}&=(Rs_n-|q_{ni}|)_+,&
 p_{ni}&=\mathbf1_{|q_{ni}|\le Rs_n},\\
 b_{ni}&=|w_i(y)^Te_n|+R\|e_n\|
                 +(\|w_i(y)\|+R)Ds_n,\\
 j_i^u&=N^{-1}\sum_{n\in C}p_{ni}b_{ni}s_n,&
 j_i^w&=N^{-1}\sum_{n\in C}d_{ni}(\|e_n\|+Ds_n),\\
 z_n&=\sum_i(\|w_i(y)\|+R)d_{ni},&
 B_C&=\Big[\sum_i((j_i^u)^2+(j_i^w)^2)\Big]^{1/2}
                  +\sqrt{\lambda_C}\bar P\|z\|_C.
\end{align*}
Every neuron is retained in these quantities.

\begin{proposition}[A static neighborhood enclosure]\label{prop:static}
For any compatible training selectors and $\|\theta-y\|\le R$,
\begin{equation}\label{eq:static}
 \left|g(\theta)^TF_C(\theta)-g(y)^TF_C(y)\right|
 \le R\left[H_g(\|F_C(y)\|+L_CR)+G_0L_C\right]
                   +(G_0+H_gR)B_C.
\end{equation}
This applies to the total derivative and separately to each cluster's
signed contribution. Cluster contributions always use the same full residual.
\end{proposition}
\begin{proof}
The Jacobian norm on cluster $C$ is at most
$\sqrt{\lambda_C}\|\theta\|$: apply Cauchy--Schwarz across neuron and
input/output parameter coordinates before taking the empirical average.
Integration along a parameter segment gives
$\|f_\theta-f_y\|_C\le\sqrt{\lambda_C}D$ and the pointwise bound
$\|f_\theta(x_n)-f_y(x_n)\|\le Ds_n$.
The same bounds hold for the fixed-mask polynomial extension.
Its field derivative is $-J_C^*J_C+H_{e,C}$, where the neuronwise
residual Hessian has norm at most
$\sqrt{\lambda_C}\|e\|_C$. Thus
$\|DF_C^0\|\le L_C$. Differentiating either potential gives
$\|Dg\|\le H_g$ on the ball. Adding and subtracting
$g(y)^TF_C^0(\theta)$ proves the first term of \eqref{eq:static}.

A changed training activation differs from its fixed-mask extension by at
most $d_{ni}$, and its selector can differ only when $p_{ni}=1$.
The pointwise residual and weight bounds give
$|w_i(\theta)^Te_\theta(x_n)|\le b_{ni}$. The direct input- and
output-gradient changes are therefore bounded by $j_i^u,j_i^w$.
The change of residual between the actual network and the fixed-mask
extension has norm at most $z_n$. Applying the fixed-mask adjoint to
this residual costs at most $\sqrt{\lambda_C}\bar P\|z\|_C$.
Consequently $\|F_C-F_C^0\|\le B_C$. Multiply by
$\|g(\theta)\|\le G_0+H_gR$ to finish. Endpoint partial selectors obey
the same bounds.
\end{proof}

\paragraph{Verified static values.}
The two centers are RK4 reference states at $3.4$ and $3.9$, generated with
step $1/5000$, the original canonical dataset and the exact archived initial
parameter values. The cohort is the original offline strong cohort of
$71$ neurons. The bounds below hold on balls of radius $3.5\cdot10^{-5}$.
Entries in the table are outward rounded and multiplied by $10^3$.
\begin{center}
\begin{tabular}{llrr}
\toprule
Statistic & contribution & center at $3.4$ & center at $3.9$\\
\midrule
$Q_G$ & total & $[-1.205,-0.894]$ & $[0.677,0.982]$\\
$Q_G$ & strong cluster & $[-2.479,-2.198]$ & $[-1.224,-0.944]$\\
$Q_G$ & weak cluster & $[1.211,1.367]$ & $[1.835,1.992]$\\
$\dot E$ & total & $[-1.439,-1.156]$ & $[0.449,0.728]$\\
$\dot E$ & strong cluster & $[-2.522,-2.266]$ & $[-1.234,-0.978]$\\
$\dot E$ & weak cluster & $[1.025,1.168]$ & $[1.622,1.766]$\\
\bottomrule
\end{tabular}
\end{center}
Thus weak-cluster forcing opposes the strong-cluster improvement at both
states and dominates it at the later state. This is a signed decomposition
of the actual field, not a retraining ablation or an assumed positive source.
The full-error computation includes all probe-active neurons directly;
it does not infer its sign from a small forward reduction residual.

On both balls the cohort lies in the positive strong cone: input angles
are below $15^\circ$, output angles below $20^\circ$. Its ungated
$m_{11}$ is greater than $0.9322$ and $0.9340$, respectively.
These are coarse learned-block geometry statements, not fine locking or
an initialization-to-specialization theorem. The full-field speed is at most
$0.385$ and $0.309$, respectively.

The accompanying script evaluates \eqref{eq:static} with outward
midpoint-radius bounds, dot-product rounding bounds, and interval
trigonometry. All center training masks are resolved by intervals; possible
changes anywhere in each ball are charged by $B_C$. The probe masks stay
fixed because their center margins exceed $3.7575\cdot10^{-5}$ and
$5.6601\cdot10^{-5}$. Exact-rational corner tests independently check
$553$ arithmetic enclosures. The arithmetic contract is IEEE
round-to-nearest with standard dot accumulation and no unsafe reassociation;
these tests are not a proof of initialized reachability.
